\documentclass[conference,a4paper]{IEEEtran}

\usepackage{amsmath,amssymb,graphicx,cite,booktabs,balance}
\usepackage[hidelinks]{hyperref}
\usepackage{url}
\usepackage{xcolor}

\hypersetup{
  pdftitle={X-Fusion: A Tri-Modal Transformer System for Multi-Label Chest X-ray Screening},
  pdfauthor={Sidhardh Ramineni, Vura Ramesh, Saharsh Batte, P. Deepesh, R. Vignesh}
}

\begin{document}

\title{X-Fusion: A Tri-Modal Transformer System for Multi-Label Chest X-ray Screening}

\author{\IEEEauthorblockN{Sidhardh Ramineni, Vura Ramesh, Saharsh Batte, P. Deepesh, R. Vignesh}
\IEEEauthorblockA{Department of Computer Science and Engineering\\
Keshav Memorial Institute of Technology\\
Hyderabad, Telangana, India\\
Emails: sidhardhramini2007@gmail.com, vuraramesh01@gmail.com, \newline saharshbatte356@gmail.com, deepeshp113@gmail.com, rvignesh72007@gmail.com}}

\maketitle

\begin{abstract}
Chest radiographs are commonly used for the screening and diagnosis of thoracic diseases because they are inexpensive, fast to acquire, and widely available. However, image-only analysis can miss clinically relevant context and is often limited by overlapping anatomical structures, subtle abnormalities, and multiple concurrent findings in a single scan. To address these challenges, we present X-Fusion, a tri-modal research prototype that combines a chest X-ray image, a clinical note, and structured patient metadata for multi-label disease screening. The proposed system integrates a pretrained biomedical visual encoder, a clinical text encoder, and a metadata encoder into a compact Transformer fusion module that outputs independent disease scores for multiple abnormalities. The framework is designed as an explainable and reproducible screening prototype that supports multi-label inference, comparison of modality contributions, and interpretability through attention-based explanations. In this work, we describe the end-to-end architecture, the training objective, and the implementation workflow, while emphasizing that reliable clinical validation requires independent patient data and externally verified evaluation metrics.
\end{abstract}

\begin{IEEEkeywords}
Chest radiography, multimodal learning, medical image analysis, Vision Transformer, clinical metadata, multi-label classification.
\end{IEEEkeywords}

\section{Introduction}

Chest radiography remains one of the most widely used medical imaging modalities for assessing thoracic diseases. It is accessible, inexpensive, and well suited for initial screening of abnormalities such as pneumonia, cardiomegaly, pleural effusion, and pneumothorax. Nevertheless, diagnosis from chest radiographs is difficult because several diseases can present with overlapping visual patterns, and a single image may contain multiple abnormalities simultaneously. For this reason, diagnostic systems often benefit from additional clinical context beyond the image itself.

In recent years, deep learning methods have demonstrated strong performance in medical image recognition tasks, especially when adapted from large-scale pretrained models. Vision Transformers, biomedical image encoders, and clinical language models have all been applied to radiology problems with promising outcomes. However, many systems still focus on image-only classification and do not fully exploit the complementary information that accompanies a patient encounter, such as clinical notes and structured vitals. This motivates the development of multimodal architectures that combine visual and non-visual signals in a principled way.

This paper presents X-Fusion, a tri-modal Transformer-based framework for multi-label chest X-ray screening. The system takes three inputs: a chest radiograph, a clinical note, and metadata such as age, sex, temperature, heart rate, respiratory rate, and oxygen saturation. Each modality is encoded separately before being fused through a compact Transformer module, which produces a set of disease scores. The framework is designed to be interpretable and reproducible, while also clarifying the difference between model optimization and clinically validated predictive performance.

The main objective of this work is not only to design a high-performing model, but also to document the engineering workflow behind a multimodal screening prototype. We describe the system architecture, the data flow, and the training strategy, while noting that valid clinical claims require independent patient data, careful evaluation protocols, and external validation against expert-labeled outcomes.

\section{Related Work}

Chest X-ray classification has been a major area of research in medical imaging. The NIH ChestX-ray14 dataset introduced a widely adopted benchmark for multi-label classification on thoracic abnormalities and remains a common reference point in the literature \cite{wang2017}. CheXpert further expanded the evaluation landscape by providing chest radiographs with uncertainty-aware labels and radiologist comparisons \cite{irvin2019}. These datasets have helped establish standard baselines for disease recognition, but they also highlight the importance of careful interpretation when labels are derived from weak supervision or auxiliary observations.

Vision Transformers have become an important architectural choice for image-based medical analysis. The original Vision Transformer demonstrated that patch embeddings and self-attention can achieve strong performance on large-scale visual recognition tasks \cite{dosovitskiy2021}. BiomedCLIP extends this idea into the biomedical domain by learning representations from scientific image-text pairs, making it well suited for medical image understanding \cite{biomedclip2023}. In parallel, clinical language models such as CXR-BERT have been developed to capture text semantics relevant to radiology and patient history \cite{cxrbert2022}.

In multi-label learning, class imbalance and varying prevalence across findings are significant challenges. Asymmetric Loss (ASL) was proposed to address imbalance by applying different treatment to positive and negative labels in multi-label learning \cite{ridnik2021}. Explainability methods such as SHAP and attention-based interpretability techniques also provide useful tools for understanding model decisions, though they should be interpreted carefully in clinical settings \cite{lundberg2017,chefer2021}. These works motivate the design of X-Fusion as a multimodal, interpretable screening framework rather than a clinically validated diagnostic model.

\section{System Overview}

The proposed system adopts a multimodal architecture in which each input stream is processed independently before fusion. The image branch uses a pretrained biomedical encoder to produce a pooled visual representation. The text branch takes a clinical note as input and generates a semantic representation using a radiology-oriented language model. The metadata branch maps a six-dimensional patient context vector into a dense embedding through a small neural network.

The resulting modality-specific vectors are concatenated with a learned classification token and passed to a lightweight Transformer encoder. The fusion module captures interactions among the three representations and produces a shared contextual representation for the final prediction head. A linear classifier followed by sigmoid activations generates one score per disease label. This design allows the system to model both individual modality characteristics and cross-modality dependencies while preserving a compact architecture suitable for research prototyping and end-to-end deployment.

\section{Methodology}

Let $I$ denote the input chest radiograph, $T$ the clinical note, and $m$ the structured metadata vector. The image encoder produces a feature vector $z_v$, the text encoder produces $z_t$, and the metadata encoder produces $z_m$:

\begin{equation}
z_v = f_v(I) \in \mathbb{R}^{512},
\end{equation}

\begin{equation}
z_t = W_t h_t(T) + b_t \in \mathbb{R}^{512},
\end{equation}

\begin{equation}
z_m = f_m(m) \in \mathbb{R}^{512}.
\end{equation}

These three vectors are then combined with a learned classification token $c$ to form a fusion input sequence $S = [c, z_v, z_t, z_m]$. A two-layer Transformer encoder processes this sequence and produces a contextual representation for classification. The final classifier outputs $K$ logits, where $K$ corresponds to the number of disease labels. Independent sigmoid functions convert the logits to probabilities or scores, which are then used for multi-label screening.

This formulation keeps the architecture modular: the visual, textual, and metadata pathways can be studied independently, while the fusion layer captures cross-modal interactions. The system is therefore suitable for both research analysis and practical deployment experiments in a clinical decision-support workflow.

\section{Implementation Considerations}

A practical deployment of X-Fusion requires careful treatment of data preprocessing, feature alignment, and model inference. The image input is normalized to a fixed shape and processed by the vision encoder. The note input is tokenized and truncated or padded to a consistent length before being fed to the text encoder. Metadata fields are standardized so that numerical values from age, sex, temperature, heart rate, respiration rate, and oxygen saturation are comparable during training.

The fusion layer integrates these features into a unified representation that is then mapped to disease logits. Because the task is multi-label, each finding is treated independently, and the final output is a vector of class-wise scores rather than a single class label. This design is aligned with the real-world setting of chest X-ray screening, where a patient may have multiple concurrent findings.

The software workflow also supports interpretability. Attention visualization and feature attribution can be used to inspect which modality or token contributed most strongly to a given prediction. These explanations are useful for model debugging and for communicating system behavior, though they do not replace expert clinical validation.

\section{Experimental Perspective and Limitations}

The development of the X-Fusion prototype highlights the difference between technical optimization and clinical reliability. During training, the loss curve can show consistent improvement, indicating that the model is learning to reduce the objective function on the available dataset. However, such optimization evidence alone does not imply diagnostic performance in real clinical settings. This is especially important when the data pipeline generates synthetic notes or structured values using the same target labels that supervise the model.

In such cases, target leakage can inflate apparent predictive performance because the input features are partially derived from the labels being predicted. For a clinically meaningful evaluation, inputs should be independent of the diagnostic targets, and the model should be assessed on patient-disjoint validation and test sets with external labels from expert review. These safeguards are essential before making any claim of real-world disease screening capability.

\section{Conclusion}

This paper presents X-Fusion as a tri-modal research prototype for multi-label chest X-ray screening. By combining image, text, and metadata information in a single Transformer-based architecture, the system demonstrates a practical and interpretable multimodal workflow for disease inference. The design is modular, reproducible, and suitable for further experimentation in medical AI research.

At the same time, the implementation emphasizes a critical point: model optimization and software functionality are not equivalent to clinical validity. Future work should focus on independent clinical data acquisition, careful separation of training and evaluation pipelines, and rigorous external validation. Only then can multimodal chest X-ray models be confidently evaluated for actual screening and diagnostic utility.

\begin{thebibliography}{99}

\bibitem{wang2017}X. Wang et al., ``ChestX-ray8: Hospital-scale chest X-ray database and benchmarks on weakly supervised classification and localization of common thoracic diseases,'' in \emph{Proc. IEEE Conf. Comput. Vis. Pattern Recognit.}, 2017, pp. 2097--2106.

\bibitem{irvin2019}J. Irvin et al., ``CheXpert: A large chest radiograph dataset with uncertainty labels and expert comparison,'' in \emph{Proc. AAAI Conf. Artif. Intell.}, vol. 33, 2019, pp. 590--597.

\bibitem{dosovitskiy2021}A. Dosovitskiy et al., ``An image is worth 16x16 words: Transformers for image recognition at scale,'' in \emph{Int. Conf. Learn. Representations}, 2021.

\bibitem{biomedclip2023}S. Zhang et al., ``BiomedCLIP: A multimodal biomedical foundation model pretrained from fifteen million scientific image-text pairs,'' arXiv:2303.00915, 2023.

\bibitem{cxrbert2022}B. Boecking et al., ``Making the most of text semantics to improve biomedical vision-language processing,'' in \emph{European Conf. Comput. Vis.}, 2022.

\bibitem{ridnik2021}T. Ridnik et al., ``Asymmetric loss for multi-label classification,'' in \emph{Proc. IEEE/CVF Int. Conf. Comput. Vis.}, 2021, pp. 82--91.

\bibitem{lundberg2017}S. M. Lundberg and S.-I. Lee, ``A unified approach to interpreting model predictions,'' in \emph{Advances in Neural Information Processing Systems}, vol. 30, 2017.

\bibitem{chefer2021}H. Chefer, S. Gur, and L. Wolf, ``Transformer interpretability beyond attention visualization,'' in \emph{Proc. IEEE/CVF Conf. Comput. Vis. Pattern Recognit.}, 2021, pp. 782--791.

\end{thebibliography}

\end{document}
